
% ==============================================
% preamble.tex
% ユーザー定義のコマンド、環境、レイアウト設定などを記述
% ==============================================

% ---- レイアウト設定 (geometry) ----
\usepackage{geometry}
\geometry{top=25mm,bottom=25mm,left=25mm,right=25mm}

% ---- 数式番号を章ごとに (1.1) のようにする ----
\numberwithin{equation}{section}

% ---- mathpad: 数式ブロックにだけ余白を付ける環境 ----
\newtcolorbox{mathpad}{%
  enhanced,
  frame hidden,
  boxrule=0pt,
  colback=white,
  opacityback=0,
  left=6mm,right=6mm,
  top=2mm,bottom=2mm,
  boxsep=0pt,
  before skip=2pt,
  after skip=2pt,
  before upper={\setlength{\jot}{5pt}},
}

% ---- pstep: (1) みたいな手順表示 ----
\newlength{\pstepindent}
\setlength{\pstepindent}{3.2em}

\NewDocumentEnvironment{pstep}{m}{%
  \par\noindent
  \makebox[\pstepindent][l]{\textbf{(#1)}}%
  \par\noindent\hspace*{\pstepindent}%
  \begin{minipage}[t]{\dimexpr\linewidth-\pstepindent\relax}%
  \setlength{\abovedisplayskip}{2pt}%
  \setlength{\belowdisplayskip}{2pt}%
  \setlength{\abovedisplayshortskip}{1pt}%
  \setlength{\belowdisplayshortskip}{1pt}%
  \ignorespaces}
  {%
  \end{minipage}\par
}

% ---- \xoverline{...}: 少し長めの上線（複素共役などで使う） ----
\makeatletter
\newsavebox{\xoverline@boxA}
\newsavebox{\xoverline@boxB}
\newcommand*\xoverline[2][0.75]{%
  \sbox{\xoverline@boxA}{$\m@th#2$}%
  \setbox\xoverline@boxB\null%
  \ht\xoverline@boxB=\ht\xoverline@boxA%
  \dp\xoverline@boxB=\dp\xoverline@boxA%
  \wd\xoverline@boxB=#1\wd\xoverline@boxA%
  \sbox{\xoverline@boxB}{$\m@th\overline{\copy\xoverline@boxB}$}%
  \setlength\@tempdima{\wd\xoverline@boxA}%
  \addtolength\@tempdima{-\wd\xoverline@boxB}%
  \ifdim\wd\xoverline@boxB<\wd\xoverline@boxA%
    \rlap{\hskip 0.5\@tempdima\usebox{\xoverline@boxB}}%
    \usebox{\xoverline@boxA}%
  \else
    \usebox{\xoverline@boxB}%
  \fi
}
\makeatother

% ---- tcolorbox 互換: \tcbbreak が無い環境用 ----
\makeatletter
\providecommand{\tcbbreak}[1][]{%
  \par\pagebreak[2]\par
}
\makeatother

% ---- その他 演算子定義 ----
\DeclareMathOperator{\Arg}{Arg}

\newcommand{\df}{\mathrel{\mathrm{df}}}
\newcommand{\Int}{\mathrm{I}}
\newcommand{\Up}{\mathrm{Up}}
\newcommand{\Un}{\mathrm{Un}}
\newcommand{\true}{\mathrm{真}}
\newcommand{\false}{\mathrm{偽}}
\newcommand{\logicnote}[1]{\textcolor{TBred}{#1}}
